\begin{frame}[fragile]{Como o template usa Lua}
  \begin{columns}[T,onlytextwidth]
    \begin{column}{0.45\textwidth}
      \begin{destaque}{Um módulo, três usos}
        \begin{itemize}
          \item Constantes e unidades
          \item Tabelas e listas de exercícios
          \item CSV $\rightarrow$ gráficos \texttt{pgfplots}
        \end{itemize}
      \end{destaque}
      \begin{destaque}[fisicaviolet]{Compilar}
        \small
        \texttt{latexmk -lualatex main.tex}\\[2pt]
        ou simplesmente \texttt{make}
      \end{destaque}
      {\scriptsize\color{fisicamuted}%
       No texto: \texttt{\textbackslash cte\{hbar\}} e
       \texttt{\textbackslash energia\{0\}\{250\}}.}
    \end{column}
    \begin{column}{0.52\textwidth}
      \begin{lstlisting}[style=lua]
-- lua/fisica.lua (trecho)
function M.constante(chave)
  local c = constantes.obter(chave)
  tex.sprint(string.format(
    "$%s = \\qty{%s}{%s}$",
    c.simbolo,
    c.valor,
    c.unidade))
end
      \end{lstlisting}
    \end{column}
  \end{columns}
\end{frame}
